% Adapted from introduction_EN_v5.md; research claims retained.
\section{Introduction}
\label{sec:introduction}

Distributed cyber-physical computation requires nodes running on independent clocks to coordinate their computations and communications. For applications that depend on reproducible execution, the relationships between communicated events must remain well defined as physical timing varies. This creates a fundamental implementation challenge: a coherent logical execution must be realized using oscillators, communication links and storage resources whose physical behaviour is neither identical nor unlimited.

Logical synchrony provides a framework for addressing this challenge without distributing a global physical clock. In bittide, continuous frame exchange and elastic receiving buffers connect the physical operation of neighbouring nodes to their logical execution. Each node adjusts its oscillator frequency using locally observed buffer occupancies. These buffers consequently serve two roles: they provide the feedback signal for synchronization and absorb the imbalance between incoming frames and local consumption. The storage available to accommodate that imbalance is finite~\cite{lall2024logicalSynchrony}.

The resulting control problem concerns the entire trajectory towards an operating point. Initial oscillator differences and changes in frame-delivery time can consume buffer headroom before the network settles or while it responds to a disturbance. A small eventual frequency disagreement does not by itself establish that the intervening buffer trajectory remains admissible. The objective is therefore to use local occupancy feedback to reduce frequency disagreement and control deviations from compatible occupancy targets, while respecting finite buffer capacity and oscillator-correction limits. This formulation makes transient buffer excursion, convergence and control activity joint criteria for evaluating a controller.

Existing bittide research establishes important foundations for this problem. PI performance analysis~\cite{lall2022resistance} relates closed-loop behaviour to controller gains and network resistance distances. Fluid occupancy analysis~\cite{lall2024bufferOccupancy} characterizes steady buffer levels, and reframing methods~\cite{lall2023bufferCentering} address buffer centering. Our earlier study, \emph{Stabilizing the bittide by Tuning Controllers Under Topology and Latency Constraints}, examined PI gain selection under topology and fixed-delay constraints. Building on these results, the question addressed here is how a local nonlinear Fuzzy PI law can reshape transient resource use under comparable scales, memory and update conventions. In particular, the connection between the nonlinear rule surface, a provable buffer-excursion benefit and behaviour under delivery perturbations requires explicit analysis and controlled comparison.

We propose a decentralized occupancy-based Fuzzy PI controller for bittide. Each node computes an oscillator-correction increment from its receiving-buffer error and the update-to-update change in that error. In the central input region studied here, the proposed surface strengthens its calibrated linear counterpart for same-sign inputs and attenuates its magnitude for opposing-sign inputs. This direction-dependent action provides flexibility at fixed calibrated PI gain coordinates. It motivates investigating whether the nonlinear correction can reduce transient buffer demands while retaining convergence, and determining how its effects on frequency disagreement and control activity depend on the operating conditions.

This paper presents the complete controller, establishes its con\-tinuous-clock realization and combines nonlinear analysis with evaluation to investigate that question. We derive the exact central-region map and its scale correspondence with native PI, and verify implementation equivalence when the nonlinear map is replaced by its calibrated linear counterpart. For a synchronous, zero-delay fluid reference, we establish an invariant input region and a common Lyapunov certificate. An outward interval enclosure of the nonlinear trajectory, followed by a contracting infinite-tail bound, provides an all-time buffer-excursion guarantee. For one prescribed three-node initialization at matched gain coordinates, the Fuzzy excursion is certified below 17.315 fluid frames, whereas PI attains an exact peak of approximately 23.403701 frames. This establishes a strict buffer-headroom benefit for that reference and initialization.

We separately evaluate the controllers in the continuous-phase Callisto model, with quantized buffer events, local-clock polling and control-execution delay. Controller updates change the slope of continuous phase trajectories; they do not replace the plant with a discrete firing-count model. Shared frame-delivery stress profiles include a latency burst after nominal centering and bounded random-latency traces. Over the declared measurement window beginning at burst onset, Fuzzy PI has lower integrated frequency-disagreement amplitude and command total variation, alongside a higher peak buffer excursion. In both random-latency regimes, the matched moderate PI baseline achieves lower disagreement RMS over the declared disturbance windows. These experiments characterize the performance tradeoffs accompanying nonlinear response shaping and distinguish the reference-model guarantee from its practical applicability under the declared execution assumptions.

% The contributions of this paper are:

% \begin{enumerate}
%   \item \textbf{Fuzzy PI control for bittide.} A proposed decentralized occupancy-feedback controller with a fully specified inference law, initialization and continuous-clock actuation, together with PI calibration and exact nonlinear analysis that support controlled comparison.
%   \item \textbf{Convergence and transient buffer margins.} An invariant-domain and Lyapunov analysis for the specified fluid reference, with verified correction and positive-frequency bounds, and an all-time comparison proving a strict headroom benefit for a prescribed initialization at matched gain coordinates.
%   \item \textbf{Evaluation under delivery perturbations.} A continuous-clock Callisto study using shared traces, clean counterfactuals, complete-event buffer measurements and independent paired random seeds to characterize frequency, buffer-margin and control-activity tradeoffs.
% \end{enumerate}
